\documentclass[aspectratio=169, 10pt]{beamer}

% --- Configuración de Idioma y Codificación ---
\usepackage[utf8]{inputenc}
\usepackage[T1]{fontenc}
\usepackage[spanish, es-tabla]{babel}

% --- Paquetes de Estilo y Gráficos ---
\usepackage{amsmath, amssymb}
\usepackage{booktabs}
\usepackage{listings}
\usepackage{xcolor}
\usepackage{hyperref}
\usepackage{tikz}

% --- Tema de la Presentación ---
\usetheme{Madrid}
\usecolortheme{whale}

% Ajuste de pie de página limpio
\setbeamertemplate{navigation symbols}{}

% --- Definición de Colores para Código C++ ---
\definecolor{codegreen}{rgb}{0.1, 0.55, 0.1}
\definecolor{codegray}{rgb}{0.45, 0.45, 0.45}
\definecolor{codepurple}{rgb}{0.5, 0.0, 0.7}
\definecolor{codeblue}{rgb}{0.0, 0.25, 0.75}
\definecolor{codebg}{rgb}{0.97, 0.97, 0.98}

% --- Configuración del Entorno de Código Listings ---
\lstdefinestyle{cppstyle}{
    backgroundcolor=\color{codebg},   
    commentstyle=\color{codegreen}\itshape,
    keywordstyle=\color{codeblue}\bfseries,
    numberstyle=\tiny\color{codegray},
    stringstyle=\color{codepurple},
    basicstyle=\ttfamily\scriptsize,
    breakatwhitespace=false,         
    breaklines=true,                 
    captionpos=b,                    
    keepspaces=true,                 
    numbers=left,                    
    numbersep=6pt,                  
    showspaces=false,                
    showstringspaces=false,
    showtabs=false,                  
    tabsize=2,
    frame=single,
    rulecolor=\color{codegray!30},
    morekeywords={UCLASS, GENERATED_BODY, UPROPERTY, UFUNCTION, AActor, Super, int32, FVector, FRotator}
}
\lstset{style=cppstyle}

% --- Metadatos de la Presentación ---
\title[POO en Unreal Engine - Parte 1]{Programación Orientada a Objetos en C++}
\subtitle{Mini-Proyecto Dado 3D: Parte 1 — Creación del Actor Vacío}
\author[SIS457]{Docente: SIS457 / Computación Gráfica y Videojuegos}
\institute[Universidad]{Carrera de Ingeniería de Sistemas}
\date{\today}

% =========================================================================
\begin{document}

% --- Diapositiva 1: Carátula ---
\begin{frame}
    \titlepage
\end{frame}

% --- Diapositiva 2: Hoja de Ruta Global ---
\begin{frame}{Mini-Proyecto: Hoja de Ruta de 11 Pasos}
    \begin{columns}[T]
        \begin{column}{0.48\textwidth}
            \textbf{Bloque 1: Estructura y Visualización}
            \begin{enumerate}
                \item \textbf{\color{blue}Paso 1: Actor Vacío (ADice)} $\leftarrow$ \textit{Hoy}
                \item Paso 2: Root y Componente de Malla
                \item Paso 3: Malla de Dado personalizada
                \item Paso 4: Atributo Value (1..6)
                \item Paso 5: Función ShowDiceFace
                \item Paso 6: Prueba reactiva en Tick
            \end{enumerate}
        \end{column}
        \begin{column}{0.48\textwidth}
            \textbf{Bloque 2: Interactividad y Física}
            \begin{enumerate}
                \setcounter{enumi}{6}
                \item Paso 7: Evento Click y Debug
                \item Paso 8: Cambio incremental de caras
                \item Paso 9: Cinemática de Salto (Fuerza/Gravedad)
                \item Paso 10: Rotación angular 3D en el aire
                \item Paso 11: Detección con Producto Punto
            \end{enumerate}
        \end{column}
    \end{columns}
\end{frame}

% --- Diapositiva 3: Objetivos del Paso 1 ---
\begin{frame}{Objetivos de la Sesión (Paso 1)}
    \begin{block}{¿Qué aprenderemos hoy?}
        \begin{itemize}
            \item Comprender qué es un \textbf{AActor} dentro de la jerarquía de Unreal Engine.
            \item Conectar conceptos de POO tradicional (C\#, Delphi, Unity) con C++ en Unreal.
            \item Aprender el rol del precompilador \textbf{UHT (Unreal Header Tool)}: macros \texttt{UCLASS()} y \texttt{GENERATED\_BODY()}.
            \item Analizar las etapas del ciclo de vida básico: \textbf{Constructor}, \textbf{BeginPlay()} y \textbf{Tick()}.
            \item Compilar y comprobar la instanciación en el \textbf{World Outliner}.
        \end{itemize}
    \end{block}
\end{frame}

% --- Diapositiva 4: Concepto de AActor ---
\begin{frame}{Concepto Fundamental: ¿Qué es un AActor?}
    \begin{itemize}
        \item En Unreal Engine, \textbf{AActor} es la clase base para cualquier objeto que pueda colocarse o generarse (spawnearse) en un nivel 3D.
        \item Hereda de \texttt{UObject} (la raíz de todos los objetos con recolección de basura).
    \end{itemize}

    \vspace{0.3cm}
    \begin{table}[]
        \centering
        \footnotesize
        \begin{tabular}{@{}lll@{}}
            \toprule
            \textbf{Entorno} & \textbf{Concepto Equivalente} & \textbf{Características} \\ \midrule
            \textbf{Unreal Engine} & \texttt{AActor} & Posee Transform (posición/rotación) si tiene Root. \\
            \textbf{Unity} & \texttt{GameObject} & Contenedor de componentes. \\
            \textbf{C\# / WinForms} & \texttt{Control} & Elemento visible en un Formulario. \\
            \textbf{Delphi} & \texttt{TWinControl} & Componente con handle visual y coordenadas. \\ \bottomrule
        \end{tabular}
    \end{table}

    \begin{alertblock}{Regla de Nomenclatura Epic}
        Las clases que heredan de \texttt{AActor} llevan obligatoriamente el prefijo \textbf{A} (ejemplo: \texttt{ADice}).
    \end{alertblock}
\end{frame}

% --- Diapositiva 5: El Sistema de Reflexión (UHT) ---
\begin{frame}{El Sistema de Reflexión de Unreal (UHT)}
    \begin{itemize}
        \item C++ estándar no tiene reflexión nativa (inspección de tipos en tiempo de ejecución).
        \item Unreal utiliza el \textbf{Unreal Header Tool (UHT)} para generar metadatos automáticamente.
    \end{itemize}

    \vspace{0.3cm}
    \begin{columns}[T]
        \begin{column}{0.48\textwidth}
            \begin{block}{Macro \texttt{UCLASS()}}
                Registra la clase para:
                \begin{itemize}
                    \item Garbage Collection (gestión de memoria).
                    \item Visualización e instanciación en el Editor.
                    \item Serialización y guardado.
                \end{itemize}
            \end{block}
        \end{column}
        \begin{column}{0.48\textwidth}
            \begin{block}{Macro \texttt{GENERATED\_BODY()}}
                \begin{itemize}
                    \item Obligatoria en la primera línea de la clase.
                    \item Inyecta código autogenerado necesario para el motor.
                    \item Requiere \texttt{\#include "Dice.generated.h"} como \textbf{último include}.
                \end{itemize}
            \end{block}
        \end{column}
    \end{columns}
\end{frame}

% --- Diapositiva 6: Código Dice.h ---
\begin{frame}[fragile]{Código Fuente: Archivo de Cabecera \texttt{Dice.h}}
\begin{lstlisting}[language=C++]
#pragma once

#include "CoreMinimal.h"
#include "GameFramework/Actor.h"
#include "Dice.generated.h" // SIEMPRE el ultimo include

UCLASS()
class PUZZLEDICEGAME_API ADice : public AActor
{
    GENERATED_BODY()
    
public: 
    // Constructor: inicializa valores por defecto
    ADice();

protected:
    // Se ejecuta una sola vez al iniciar el juego
    virtual void BeginPlay() override;

public: 
    // Se ejecuta en cada fotograma del motor
    virtual void Tick(float DeltaTime) override;
};
\end{lstlisting}
\end{frame}

% --- Diapositiva 7: Código Dice.cpp ---
\begin{frame}[fragile]{Código Fuente: Implementación \texttt{Dice.cpp}}
\begin{lstlisting}[language=C++]
#include "Dice.h"

ADice::ADice()
{
    // Habilita que Tick() se llame cada frame
    PrimaryActorTick.bCanEverTick = true;
}

void ADice::BeginPlay()
{
    Super::BeginPlay();
    
    // Mensaje de diagnostico en la consola de salida
    UE_LOG(LogTemp, Warning, TEXT("ADice creado exitosamente en el mundo"));
}

void ADice::Tick(float DeltaTime)
{
    Super::Tick(DeltaTime);
    // Aqui iran las actualizaciones continuas de fisica
}
\end{lstlisting}
\end{frame}

% --- Diapositiva 8: Ciclo de Vida ---
\begin{frame}{Ciclo de Vida: ¿Dónde colocar nuestro código?}
    \begin{columns}[T]
        \begin{column}{0.32\textwidth}
            \begin{block}{\textbf{Constructor}}
                \texttt{ADice()}
                \begin{itemize}
                    \item Se ejecuta al compilar, abrir el editor y crear el CDO.
                    \item \textbf{Uso:} Inicializar atributos por defecto y crear componentes (\texttt{CreateDefaultSubobject}).
                \end{itemize}
            \end{block}
        \end{column}
        \begin{column}{0.32\textwidth}
            \begin{block}{\textbf{BeginPlay}}
                \texttt{BeginPlay()}
                \begin{itemize}
                    \item Se ejecuta al presionar \textit{Play} o cuando el Actor nace en la escena.
                    \item \textbf{Uso:} Suscribirse a eventos, buscar otros actores y preparar el estado inicial.
                \end{itemize}
            \end{block}
        \end{column}
        \begin{column}{0.32\textwidth}
            \begin{block}{\textbf{Tick}}
                \texttt{Tick(DeltaTime)}
                \begin{itemize}
                    \item Se ejecuta 60+ veces por segundo.
                    \item \textbf{Uso:} Traslación continua, cálculo de gravedad e interpolaciones físicas.
                \end{itemize}
            \end{block}
        \end{column}
    \end{columns}
\end{frame}

% --- Diapositiva 9: Práctica en el Editor ---
\begin{frame}{Verificación Práctica en Unreal Editor}
    \begin{enumerate}
        \item \textbf{Compilar:} Presionar \texttt{Ctrl + Shift + B} en Visual Studio o el botón \textbf{Live Coding} en Unreal.
        \item \textbf{Localizar en Content Browser:} Ir a \texttt{C++ Classes} $\rightarrow$ \texttt{PuzzleDiceGame} $\rightarrow$ \texttt{Dice}.
        \item \textbf{Instanciar:} Arrastrar \texttt{Dice} hacia el Viewport 3D del nivel.
        \item \textbf{Comprobar World Outliner:} Verificar que aparece \texttt{Dice1}.
    \end{enumerate}

    \vspace{0.2cm}
    \begin{alertblock}{Pregunta para la clase}
        ¿Por qué el actor aparece en la lista del Outliner pero \textbf{no se ve ninguna forma visual} en la pantalla 3D?
        \newline
        $\rightarrow$ \textit{Porque un \texttt{AActor} es solo un contenedor lógico; requiere \textbf{Componentes} para tener forma visible.}
    \end{alertblock}
\end{frame}

% --- Diapositiva 10: Control de Versiones Git ---
\begin{frame}[fragile]{Registro del Hito en Git}
    Para que los alumnos puedan descargar y verificar este estado del proyecto:

    \vspace{0.2cm}
    \begin{lstlisting}[language=bash]
# 1. Comprobar archivos modificados
git status

# 2. Agregar archivos de cabecera e implementacion
git add Source/PuzzleDiceGame/Dice.h Source/PuzzleDiceGame/Dice.cpp

# 3. Crear el commit descriptivo
git commit -m "feat(dice): Paso 01 - Creacion del actor vacio ADice"

# 4. Crear etiqueta de version para la clase
git tag step-01-empty-actor
    \end{lstlisting}
\end{frame}

% --- Diapositiva 11: Conclusión y Próximo Paso ---
\begin{frame}{Resumen y Siguiente Sesión}
    \begin{block}{Lo que hemos logrado hoy:}
        \checkmark Estructura base C++ de la entidad \texttt{ADice} registrada en el motor.\\
        \checkmark Comprensión del sistema de reflexión \texttt{UCLASS} y \texttt{GENERATED\_BODY}.\\
        \checkmark Verificación de instanciación en el nivel.
    \end{block}

    \vspace{0.4cm}
    \begin{exampleblock}{Próximo Paso: Paso 02 — Composición y Geometría}
        Aprenderemos el principio de \textbf{Composición sobre Herencia}:
        \begin{itemize}
            \item Crearemos un \texttt{USceneComponent* DummyRoot}.
            \item Adheriremos un componente de malla estática \texttt{UStaticMeshComponent* BlockMesh}.
            \item Veremos el cubo 3D renderizado físicamente en el nivel.
        \end{itemize}
    \end{exampleblock}
\end{frame}

\end{document}
